Good figure captions help paper readers understand complex scientific figures.
Unfortunately, even published papers often have poorly written captions.
Automatic caption generation could aid paper writers by providing good starting captions that can be refined for better quality.
Prior work often treated figure caption generation as a vision-to-language task.
In this paper, we show that it can be more effectively tackled as a \textbf{text summarization} task in scientific documents.
%PEGASUS, an abstractive summarization model~\cite{zhang2020pegasus}
%We fine-tuned PEGASUS, an abstractive summarization model, 
We fine-tuned PEGASUS, a pre-trained abstractive summarization model, to specifically summarize figure-referencing paragraphs (\eg, ``Figure 3 shows...'') into figure captions.
Experiments on large-scale arXiv figures show that our method outperforms prior vision methods in both automatic and human evaluations.
%\kenneth{TODO: Update the numbers} \cy{Do you mean test set or all?}
%\kenneth{I took out the numbers. Is this description accurate?}
We further conducted an in-depth investigation focused on two key challenges: {\em (i)} the common presence of low-quality author-written captions and {\em (ii)} the lack of clear standards for good captions.
Our code and data are available at: \url{https://github.com/Crowd-AI-Lab/Generating-Figure-Captions-as-a-Text-Summarization-Task}.



%------------------------------ dead kitten ---------------------------------------

\begin{comment}

\bigskip

Our experiments with 1,XXX,XXX arXiv figures and captions show that summarizing paragraphs that mention figures (\eg, ``Figure 4 shows...'') produces better captions according to both automatic and human evaluations.



Prior work often treated figure caption generation as a vision-to-language task.
In this paper, we show that it can be effectively tackled as a \textbf{text summarization} task in scientific documents.


Good figure captions help paper readers understand complex scientific figures.
Unfortunately, poorly-written captions occur frequently, even in published papers.
Automatic caption generation can aid paper writers by providing decent initial captions; writers can then revise these machine-generated captions to craft higher-quality ones.
Prior work often treated figure captioning for scholarly articles as a vision-to-language task. 
In this paper, we show that it can be effectively tackled as a \textbf{text summarization} task.







%\textbf{\LARGE
%\textcolor{red}{To change the title or to find your commenting command: see {\tt %macro.tex}!}
%}

%\kenneth{}


%The task of generating captions for scientific figures presents a significant challenge for artificial intelligence, as it requires a computational model to not only comprehend the visual content of the figure, but also to grasp the complex contextual information present in the accompanying scientific document. 
%Generating captions for scientific figures is a significant AI challenge, requiring a model to comprehend visual content and grasp complex contextual information. %\kenneth{Maybe rephrase this, as it turns out we do not REQUIRE visual features?}
Effective figure captions are crucial for clear comprehension of scientific figures, yet poor caption writing remains a common issue in scientific articles.
%
%Clear and effective figure captions play a critical role in facilitating the %However, poor caption writing remains a widespread issue in scientific articles despite its importance.
%Effective figure captions are vital for clear comprehension of scientific figures, yet poor caption writing continues to plague scientific articles.
Our study of arXiv {\tt cs.CL} papers found that 53.88\% of captions were rated as unhelpful or worse by domain experts, 
%Our study of arXiv cs.CL papers found that a substantial 53.88\% of captions were rated as unhelpful or worse by domain experts, 
showing the need for better caption generation.
Previous efforts in figure caption generation treated it as a vision task, aimed at creating a model to understand visual content and complex contextual information.
Our findings, however, demonstrate that over 75\% of figure captions' tokens align with corresponding figure-mentioning paragraphs, indicating great potential for language technology to solve this task.
%demonstrate that more than 75\% of tokens in figure captions can be aligned with the corresponding figure-mentioning paragraphs, suggesting that language technology is a more effective solution for this task.
%To address this, 
In this paper, we present a novel approach for generating figure captions in scientific documents using \textbf{text summarization} techniques. 
%In this work, we present a novel approach for generating figure captions in scientific documents using \textbf{text summarization} techniques.
Our approach extracts sentences referencing the target figure, then summarizes them into a concise caption.
%Figures play an essential role in scientific articles., but they are not always easy to understand.
%Good captions are necessary to understand the figures in scientific articles, 
%but people often write bad captions.
%Scientists use charts in papers to show experimental results; companies use charts in business reports to show financial status; journalists use charts in news articles to show trends and statistics.
%\kenneth{Can use one more or one better example here where *a chart is in a document*.}\ed{government use chart in the official documents to show the trends and disease statistics}.
%When glancing through a full-length document, the visuals, such as charts, graphics, and diagrams, help readers understand the content effectively.
%In this paper, we propose a novel approach to figure caption generation in scientific documents by treating it as a \textbf{text summarization task.} 
%Our method leverages the full text of the document to generate captions that are more informative and effective than traditional vision-to-language approaches.
%Our method utilizes the complete document text to generate captions that surpass traditional vision-to-language methods in terms of informativeness and effectiveness.
%Our approach first extracts all sentences in the document explicitly referencing the target figure (\eg, ``As shown in Figure 2...''). 
%These sentences are then summarized into a coherent and concise caption.\kenneth{Maybe mention the model we use.}
%Our approach extracts sentences referencing the target figure, then summarizes them into a concise caption.
%In our experiments on XXX,XXX real-world arXiv  papers, our method achieved superior results compared to previous approaches, despite relying solely on textual information.\kenneth{TODO TY: Update the numbers + how many often are published?}
In the experiments on real-world arXiv papers (81.2\% were published at academic conferences), 
our method, using only text data, outperformed previous approaches in both automatic and human evaluations.
%our method achieved superior results compared to previous approaches, despite relying solely on textual information.
%To move the needle in figure captioning, 
We further conducted data-driven investigations into the two core challenges: 
{\em (i)} low-quality author-written captions and 
{\em (ii)} the absence of a standard for good captions.
%\sukim{The following statement sounds like a motivating example, rather than the result showing the superiority of our method. It needs be moved to the beginning of the abstract to use it as a motivation, or be highlighted more how helpful the captions generated by our method are over the baselines. -- we further conducted in-depth investigations into the two core challenges: low-quality author-written captions and the absence of a standard for good captions.
%Our study found that 
%53.88\% of the author-written captions in arXiv {\tt cs.CL} papers were rated as unhelpful or worse by domain experts in our study, motivating the need for high-quality caption generation.}
%\kenneth{TODO CY: Insert the two most interesting finding from the quality beam study.}
We found that our models could generate improved captions for figures with original captions rated as unhelpful, and the model trained on captions with more than 30 tokens produced higher-quality captions. 
We also found that good captions often include the high-level takeaway of the figure.
% \kenneth{TODO CY: Read these to make sure they're accurate.}
%The paper concludes with a data-driven analysis of the relationships between captions, paragraphs, and figure images.
Our work proves the effectiveness of text summarization in generating figure captions for scholarly articles, outperforming prior vision-based approaches.
Our findings have practical implications for future figure captioning systems, improving scientific communication clarity.
%Our paper breaks new ground by demonstrating the superiority of text summarization in generating captions for figures.
%Our findings have the potential to advance the field of figure caption generation and improve the clarity of scientific communication.

%\kenneth{TODO TY: Update the numbers. How many of the arXiv papers were published?}
%\ed{To reduce the quality concerns about using arXiv papers compared to peer-reviewed papers. We manually verified the subset of real-world arxiV papers (399 figures in total) and found that 81.2\%(324/399) were published at academic conferences, and 51.9\%(207/399) were at ACL Anthology/IEEE/ACM.}
%\sukim{I think we need to add a few more sentences summarizing as many findings as we observed from our extensive analysis since the method part is relatively simple.}

\bigskip

To shed light on the relationship between a figure's caption and the paper paragraphs that mention the figure, we conducted a comprehensive analysis.
Our findings indicate that more than 75\% of the tokens in author-written captions can be aligned with the corresponding figure-mentioning paragraphs, reinforcing the validity of our approach.
Finally, to drive further advancements in the field of figure captioning, we conducted an additional analysis that sheds light on a major challenge in this domain - the low quality of author-written captions.

This presents a significant challenge for using these low-quality captions as training data or for evaluating the performance of captioning models.





In our experiments on XXX,XXX real-world arXiv papers\kenneth{Update the numbers}, our method outperformed previous approaches in terms of both automatic evaluation metrics and human preference ratings despite using only textual information.
To gain a deeper understanding of the relationship between a figure's caption and the paper paragraphs mentioning the figure, we conducted an extensive analysis.
%Our experiments on XXX,XXXX real-world arXiv papers\kenneth{TODO: update the numbers}
%we used these captions for figures and compare our method to existing approaches. Our 
%results 
%showed that our method outperformed previous approaches in terms of both automatic evaluation metrics and human preference ratings.
Our results show that over 75\% of the tokens in author-written captions can be aligned to the figure-mentioning paragraphs, supporting the foundation of our approach.

%and generates a coherent and concise summary for them.
%Results from experiments on the \oldDataset dataset show that our method outperformed existing approaches in terms of both automatic evaluation and human preferences.
%Our analysis further suggested that over 75\% of the caption tokens can be aligned with the figure-mentioning paragraphs, demonstrating the effectiveness of our approach. 
%Finally, we highlight the need for further research in the area, given that 53.88\% of the author-written captions in arXiv {\tt cs.CL} papers are considered unhelpful.

\kenneth{Do we change the title as reviewers suggested?}
Generating figure captions has traditionally been treated as a vision-to-language task, \ie, the model produces a caption primarily based on the image.
% However, such approaches seem not to be effective. %for readers.
%often failed to generate good-enough captions for real-world charts.
In this paper, we argue that when the full text of the document is available, the task of figure captioning can be viewed as a \textit{text-summarization task}.
%instead of a vision-to-language task.
That is, summarizing the paragraphs that mention a figure can generate high-quality captions for it.
Automatic evaluation in experiments on the \oldDataset dataset showed that summarizing the figure-mentioning paragraphs produces significantly better captions than vision-to-language approaches.
%Even simply using the figure-mentioning sentence-- without any modifications-- as the captions reached better ROUGE scores.
%For human evaluation, we recruited 
When three NLP researchers ranked captions,
%On average, 
our best-performing model's caption was preferred over the original caption 46.67\% of the time.
%Our experiments conducted on the \oldDataset dataset showed that the summarization approach drastically outperforms all other baselines in automatic and human evaluations.
Further analysis demonstrated that over 75\% of tokens in an average caption can be aligned to the tokens in the figure-mentioning paragraph.
%We also reveal a surprising number of low-quality chart captions exist in arXiv paper.\kenneth{Closing sentences need more work.}
We also reveal that 53.88\% of the captions in arXiv {\tt cs.CL} papers were considered unhelpful, motivating future research to support caption readers and writers.

The theme that emerged in our meetings seemed to be the relations between: 

\begin{itemize}
    \item figure captions, 
    \item explicit mentions of the figures, and 
    \item the paragraphs that contain the mentions. 
\end{itemize}

This paper's possible story could be about how figure captions can be approximated by summarizing or paraphrasing mentions/paragraphs, raising questions about the role of figure images in caption generation.
A side bonus could be some: in cases where human-written captions are bad, we can generate better captions using mentions.


\end{comment}
